\documentclass[a4paper,twoside,BCOR=20mm, openany]{scrreprt}
\usepackage[utf8]{inputenc}
\usepackage[T1]{fontenc}
\usepackage[ngerman]{babel}	% german hyphenation, quotes, etc
\usepackage[ngerman]{translator}
\usepackage{paralist}
\usepackage{amsfonts}
\usepackage{acronym}
\usepackage{enumerate}
\usepackage{hyperref}
\usepackage{amssymb}
\usepackage{caption}
\usepackage{multirow}
\usepackage{graphicx}
\usepackage{tabularx}
\usepackage{color}
\usepackage{wrapfig} % wrap text around figures
\usepackage{subfig} % align two pics beside each other
\usepackage[table,xcdraw,dvipsnames]{xcolor}
\usepackage{enumitem} % für custom enumerations
\usepackage{float} % for better graphics placement
\usepackage{array}
\usepackage{tikz}
\usepackage{svg} % for scalable vector-graphics (see inkscape)
\setsvg{inkscapeexe=inkscape, inkscapeopt=-z -D} % set inkscape options
\svgpath{images/} % use dir 'images/' to search for svg files

\renewcaptionname{ngerman}{\figurename}{Abb.} % replace "Abbildung" with "Abb." for all captions for graphics

\hypersetup{ 					% ‘texdoc hyperref‘ for options
	pdftitle={PSE TECO: Qualitätssicherungsbericht},
	pdfauthor={Jean Baumgarten, Oliver Liu, Patrick Ries, Erik Wessel, Thomas Frank},
	bookmarks=true,
}
\title{Qualitätssicherungsbericht: PaVoS}

%Paket laden
\usepackage[
numberedsection,
nonumberlist, %keine Seitenzahlen anzeigen
acronym,      %ein Abkürzungsverzeichnis erstellen
toc,          %Einträge im Inhaltsverzeichnis
section]      %im Inhaltsverzeichnis auf section-Ebene erscheinen
{glossaries}

% KIT layout

\definecolor{orange}{rgb}{1,0.5,0}
\definecolor{mintgreen}{RGB}{50,161,137}
\definecolor{gray}{RGB}{120,120,120}

\usepackage[color]{changebar}
\cbcolor{gray}
\changebarwidth 0.5pt

\usepackage{fancyhdr}
\pagestyle{fancy}
 \fancyhf{} %alle Kopf- und Fußzeilenfelder bereinigen 
 
 \fancypagestyle{plain}{} %Kopf- und Fußzeile auf jeder Seite	 
	\fancyhead[L]{Qualitätssicherungsbericht}
	\fancyhead[R]{\leftmark}
	\rhead{\nouppercase{\leftmark}}
	\renewcommand{\headrulewidth}{0.5pt}
	\renewcommand{\headrule}{\hbox to\headwidth{%
		\color{mintgreen}\leaders\hrule height \headrulewidth\hfill}}

\raggedbottom

	\renewcommand{\footrulewidth}{0.5pt}
	\renewcommand{\footrule}{\hbox to\headwidth{%
  		\color{mintgreen}\leaders\hrule height \headrulewidth\hfill}}				
	\fancyfoot[LE,RO]{\thepage}

\newcolumntype{L}[1]{>{\raggedright\let\newline\\\arraybackslash\hspace{0pt}}m{#1}}
\newcolumntype{C}[1]{>{\centering\let\newline\\\arraybackslash\hspace{0pt}}m{#1}}
\newcolumntype{R}[1]{>{\raggedleft\let\newline\\\arraybackslash\hspace{0pt}}m{#1}}

\def\checkmark{\tikz\fill[scale=0.4](0,.35) -- (.25,0) -- (1,.7) -- (.25,.15) -- cycle;}
\def\testGood{\cellcolor{green!25} \textcolor{Green}{\checkmark}}
\def\testOk{\cellcolor{yellow!25} \textcolor{orange}{(\checkmark)}}
\def\testBad{\cellcolor{red!25} \textcolor{Maroon}{X}}

\def\szenarioGood{\begin{table}[H] \centering \begin{tabular}{|c|} \hline \testGood \\ \hline \end{tabular} \end{table}}
\def\szenarioOk{\begin{table}[H] \centering \begin{tabular}{|c|} \hline \testOk \\ \hline \end{tabular} \end{table}}
\def\szenarioBad{\begin{table}[H] \centering \begin{tabular}{|c|} \hline \testBad \\ \hline \end{tabular} \end{table}}

\renewcommand{\familydefault}{\sfdefault}

\setcounter{tocdepth}{3}
\makeglossaries
\begin{document}
	
	\begin{titlepage}
	\begin{center}
	\includegraphics[width=0.5\linewidth]{images/TECOLogo.png}\\[0.2cm]

	\textbf{TECO Research Group}\\[0.2cm]
	Marcel Köpke\\Matthias Budde\\Till Riedel\\
	\vspace{1cm}
	
	\includegraphics[width=0.33\linewidth]{images/PaVoSLogo-Erweitert}\\[1cm]
	
	\textsc{\textbf{\LARGE Pflichtenheft}}\\
	{\small Version 0.1}\\
	
	\vspace{1cm}\hrule\vspace{0.4cm}
	\textbf{\huge Visualizing \& Mining of Geospatial Sensorstreams with Apache Kafka}\\
	\vspace{0.4cm}\hrule\vspace{1cm}
	
	{\Large Jean Baumgarten\\
	Oliver Liu\\
	Patrick Ries\\
	Erik Wessel\\
	Thomas Frank\\}

	\vspace{1cm}
	\today
	
	\end{center}
\end{titlepage}
	\tableofcontents
	\chapter{Einleitung}
Dieser Implementierungsbericht dient als Übersicht der Implementierungsphase.\newline
\linebreak
Die Implementierungsphase unseres Projekts wurde in zwei Hälften unterteilt, von denen die erste drei Wochen und die zweite vier Wochen dauerte. Bezieht man die zusätzliche Woche zur Klausurvorbereitung mit ein, summiert sich die Phasendauer auf acht Wochen.\newline
\linebreak
In der ersten Hälfte wurde hauptsächlich...\newline
\linebreak
In der zweiten Hälfte wurde schließlich...\newline
\linebreak
Nachfolgend werden Abweichungen der Implementierung vom Entwurf analysiert, eine Übersicht zu Unit-Tests gegeben und auf die implementierten Muss- und Wunschkriterien eingegangen.
	\chapter{Bugfixes}

\section{Import und Export}
\paragraph{Servlet ging nicht}
\begin{itemize}
	\item Fehlersymptom: Servlet wird von Tomcat nicht gestartet.
	\item Fehlerursache: Kafka und JSON dependency vertragen sich anscheinend nicht mit Tomcat.
	\item Fehlerbehebung: Export-Funktionalität die Kafka benötigt wurde als separate jar ausgelagert.
\end{itemize}

\paragraph{usr.home Ordner}
\begin{itemize}
	\item Fehlersymptom: Servlet und Export jar benutzen nicht die selben Dateien zur persistenten Datenspeicherung.
	\item Fehlerursache: User Home Ordner ist aus Sicht eines Servlets innerhalb des Tomcat Ordners, während eine jar einen anderen Ordner betrachtet.
	\item Fehlerbehebung: Hardcoded Verzeichnisse innerhalb des Dockers.
\end{itemize}

\paragraph{Export geht nicht}
\begin{itemize}
	\item Fehlersymptom: Export bricht nach Start sofort ab und gibt leere Datei zurück.
	\item Fehlerursache: Export wurde abgebrochen, wenn von Kafka angefragte Daten außerhalb des für den Export erwünschten Zeitraums liegen. Insbesondere auch wenn sie davor liegen.
	\item Fehlerbehebung: Es wird nun nur beendet, wenn der Zeitraum der erhaltenen Daten nach den erwünschten Daten liegt.
\end{itemize}

\paragraph{Kafka Pollzeit}
\begin{itemize}
	\item Fehlersymptom: Es wird nur ein Bruchteil der Daten exportiert.
	\item Fehlerursache: Poll Timeout Zeit an Kafka ist zu kurz.
	\item Fehlerbehebung: Poll Timeout Zeit wird erhöht.
\end{itemize}

\paragraph{Millisekunden im TimeFrame}
\begin{itemize}
	\item Fehlersymptom: Exporter stürzt bei Ausführung mit manchen Daten ab.
	\item Fehlerursache: Exporter kann nur Zeiten nach ISO 8601 ohne Millisekunden verarbeiten.
	\item Fehlerbehebung: Exporter unterstützt nun auch Millisekunden.
\end{itemize}

\paragraph{@iot.id}
\begin{itemize}
	\item Fehlersymptom: Es sind keine \texttt{@iot.id}s für Observations im Export enthalten. Stattdessen steht überall \texttt{null}, wodurch der Import dieser Daten unmöglich ist.
	\item Fehlerursache: Es wurde ein leeres Feld aus den Kafkadaten ausgelesen.
	\item Fehlerbehebung: Wechsel auf das richtige Feld.
\end{itemize}
\newpage
\section{Webinterface}
\paragraph{CORS}
\begin{itemize}
	\item Fehlersymptom: Blockieren von Http-Anfragen bei Webinterface-internem JavaScript Code
	\item Fehlerursache: Cross-Origin Resource Sharing (CORS) zwischen Webinterface und Server nicht möglich
	\item Fehlerbehebung: Konfigurieren der Server-lokalen Tomcat Konfigurationsdatei
\end{itemize}

\paragraph{Synchronisierung der Zustände}
\begin{itemize}
	\item Fehlersymptom: Fehlerhafte Synchronisierung der Zustände bei Export-, Sensortyp-, Favoriten-, und Einstellungs-Modal-Fenstern
	\item Fehlerursache: Veränderung der Auswahl/Parameter in einem Modal-Fenster werden nicht an anderes Modal-Fenster weitergegeben
	\item Fehlerbehebung: Korrekte Beobachter Struktur zwischen den Komponenten aufbauen
\end{itemize}

\paragraph{Fenstergrößenanpassung}
\begin{itemize}
	\item Fehlersymptom: Unresponsivität der Anpassung bei Veränderung der vertikalen Fenstergröße
	\item Fehlerursache: Da der Container der Karte initial eine fest definierte Höhe braucht, kann man das nicht in CSS in Prozent, also dynamisch zur Fenstergröße festlegen
	\item Fehlerbehebung: Dynamisches Styling der Höhe über JavaScript Beobachter
\end{itemize}

\newpage
\section{Bridge}
\paragraph{Main-Klasse}
\begin{itemize}
	\item Fehlersymptom: Main-Klasse mit main-Methode konnte nach kompilieren ins jar-Format bei der Ausführung nicht gefunden werden.
	\item Fehlerursache: Einige Einträge in der pom.xml fehlten.
	\item Fehlerbehebung: Hinzufügen einiger Einträge in der pom.xml.
\end{itemize}

\paragraph{Source-Ordner}
\begin{itemize}
	\item Fehlersymptom: src/main/java wurde anfangs nicht als Java Source-Ordner erkannt.
	\item Fehlerursache: Eclipse wurde falsch eingestellt.
	\item Fehlerbehebung: Eclipse wurde richtig eingestellt, Importe und Paketdefinitionen wurden entsprechend angepasst.
\end{itemize}

\paragraph{\texttt{null}-Checks}
\begin{itemize}
	\item Fehlersymptom: In einigen Methoden wurde nicht nach \texttt{null}-Objekten geprüft.
	\item Fehlerursache: Fehlende Checks.
	\item Fehlerbehebung: Checks wurden implementiert.
\end{itemize}

\paragraph{Logger}
\begin{itemize}
	\item Fehlersymptom: Konfigurationsprobleme beim Logger, nicht-einheitlicher Logger mit dem Rest von PaVoS.
	\item Fehlerursache: \texttt{java.util.logging.Logger} wurde verwendet.
	\item Fehlerbehebung: der Logger für Klassen der Bridge wurde von \texttt{java.util.logging.Logger} zu \texttt{org.slf4j.Logger} geändert.
\end{itemize}

\paragraph{\texttt{System.out} und \texttt{System.err}}
\begin{itemize}
	\item Fehlersymptom: Fehler- und Informationsnachrichten wurden über \texttt{System.out} oder \texttt{System.err} ausgegeben.
	\item Fehlerursache: Fehler- und Informationsnachrichten wurden über \texttt{System.out} oder \texttt{System.err} ausgegeben.
	\item Fehlerbehebung: Fehler- und Informationsnachrichten werden über den Logger ausgegeben.
\end{itemize}

\paragraph{MqttConsumer}
\begin{itemize}
	\item Fehlersymptom: Verbindung zu MQTT wird bei Verbindungsverlust nicht wiederaufgebaut.
	\item Fehlerursache: Die \texttt{connectionLost()} Methode wurde falsch implementiert.
	\item Fehlerbehebung: Der Verbindungsaufbau wurde nun in eine zusätzliche private Methode geschoben, die sowohl im Constructor als auch bei Verbindungsverlust ausgeführt wird.
\end{itemize}

\paragraph{Properties}
\begin{itemize}
	\item Fehlersymptom: Die Bridge startet auch bei fehlerhaften Properties, oder \texttt{PropertiesFileReader.init()} beendet die Bridge.
	\item Fehlerursache: \texttt{PropertiesFileReader.init()} führt \texttt{System.exit(-1)} bei fehlerhaften Daten aus.
	\item Fehlerbehebung: Die Main-Methode prüft ob die Initialisierung erfolgreich war und beendet dann dort die Bridge.
\end{itemize}

\paragraph{Sonarqube}
\begin{itemize}
	\item Nach der Analyse der Bridge mit Sonarqube wurden 35 von 54 Code Smells und 2 von 2 Vulnerabilities gefixed.
\end{itemize}

\section{Database}
\begin{itemize}
	\item Mehrere Code Smells wurden mit Sonarqube entdeckt und durch Erik gefixed.
\end{itemize}
\newpage
\section{Core - Verarbeitung der eingehenden Daten}
\paragraph{Nicht existierende Kafka-Topics}
\begin{itemize}
	\item Fehlersymptom: Topics wurden in Kafka nicht erstellt, welches dann teilweise zu Fehlern in den einzelnen Nutzern der Kafka Stream Funktionalität führte.
	\item Fehlerursache: Die Input Topics waren nicht in Kafka enthalten
	\item Fehlerbehebung: Ich habe eine Klasse hinzugefügt welche alle Input Topic generiert und somit garantiert dass das Kafka System alle Topics die benötigt werden enthalten sind. 
\end{itemize}

\paragraph{Übertragung von FROST zu Kafka}
\begin{itemize}
	\item Fehlersymptom: Die Daten wurden nicht von Frost auf Kafka übertragen.
	\item Fehlerursache: Es gab Null Values bei den Zeiten welche nicht erlaubt waren
	\item Fehlerbehebung: Im Importer wurde einfach gefixt, dass die Zeit nicht null sein darf.
\end{itemize}

\paragraph{Nicht erstellte Kafka-Topics}
\begin{itemize}
	\item Fehlersymptom: Die Kafka Topics von den Prozessen wurden nicht erstellt.
	\item Fehlerursache: Ich ging davon aus, dass FeatureOfInterest und Observation eine bijektive Verbindung zu einander haben.
	\item Fehlerbehebung: Ich habe das so gefixed, dass es eine surjektive Verbindung wird durch das Verwenden von einem leftJoin.
\end{itemize}

\paragraph{Fehlerhafte Nachrichten beim Export-Prozess}
\begin{itemize}
	\item Fehlersymptom: Der Export Prozess produzierte fehlerhafte Messages im Export Topic
	\item Fehlerursache: Gleicher Fehler wie bei Bug 3.
	\item Fehlerbehebung: Gleicher Fehler wie bei Bug 3.
\end{itemize}
\newpage
\section{Core - Grid und Webserver}
\paragraph{LoggerFactory - ClassNotFoundException}
\begin{itemize}
	\item Fehlersymptom: Maven auf Debian kann Pavos-Core nicht erstellen\\
	\texttt{ClassNotFoundException: LoggerFactory.java}
	\item Fehlerursache: Benötigte Dateien sind nicht im Classpath
	\item Fehlerbehebung: Maven dependency plugin fügt dependencies zum Classpath hinzu
\end{itemize}

\paragraph{Avro Manifest}
\begin{itemize}
	\item Fehlersymptom: Maven auf Debian kann Pavos-Core nicht erstellen
	Avro has no Manifest
	\item Fehlerursache: Avro.jar enthält keine Manifest-Datei
	\item Fehlerbehebung: Aus Maven dependency plugin ausschließen
\end{itemize}

\paragraph{Sensorposition}
\begin{itemize}
	\item Fehlersymptom: Erstellter GeoGrid akzeptiert korrekte Sensordaten nicht.
	\item Fehlerursache: \texttt{PointNotOnMapException}
	\item Fehlerbehebung: In der Erstellung wurde x und y als Breite und Höhe interpretiert, weshalb die Kartengröße nur 25\% der gewollten Größe betrug
\end{itemize}

\paragraph{Multigradient toString()}
\begin{itemize}
	\item Fehlersymptom: MultiGradient toString() wirft eine \texttt{NullPointerException}
	\item Fehlerursache: Farbwerte hatten einen Index-Shift um 1, weshalb der erste Farbwert übersprungen wurde und der letzte null war
	\item Fehlerbehebung: Index-Shift korrigiert und Absicherung gegen NullPointerException in ColorUtil erstellt
\end{itemize}

\paragraph{Tests mit Kafka}
\begin{itemize}
	\item Fehlersymptom: Kafka-Tests schlagen bei mvn clean install fehl
	\item Fehlerursache: JUnit tests reichen nicht aus, da Kafka zur Kommunikation benötigt wird.
	\item Fehlerbehebung: Dependency für Kafka-Streams-Test-Util für Version 1.1.0 hinzugefügt, da 1.0.1 nicht existiert
\end{itemize}

\paragraph{WebServer - automatischer Start}
\begin{itemize}
	\item Fehlersymptom: WebServer soll immer nach den KafkaProzessen gestartet werden und trotzdem soll Webserver auch selbstständig gestartet werden können.
	Beide haben main(String[] args)
	\item Fehlerursache: Zwei main Methoden können nicht ausgeführt werden, da keine zwei Main klassen existieren können.
	\item Fehlerbehebung: WebServer implementiert Runnable und ruft run() in seiner main Methode auf.
	KafkaProzesse, die Main Klasse startet WebServer auf einem Thread.
\end{itemize}

\paragraph{WebServer - kein Antwort}
\begin{itemize}
	\item Fehlersymptom: WebServer wirft Fehlermeldungen und lässt die Verbindung abstürzen.
	\item Fehlerursache: http-Anfrage ist falsch formuliert oder angefragte Komponenten existieren nicht.
	\item Fehlerbehebung: Der http-Status des WebServer wird entsprechend geändert. Meistens auf Bad-Request und wird gesendet.
\end{itemize}

\paragraph{Verbindung zur Datenbank}
\begin{itemize}
	\item Fehlersymptom: Die Datenbank kann bei der Übertragung nicht erreicht werden.
	\item Fehlerursache: Die Verbindung zur VM wurde nicht eingestellt.
	\item Fehlerbehebung: Host wurde von localhost zur VM-Host-Adresse geändert.
\end{itemize}

\paragraph{Regex - Formatüberprüfung}
\begin{itemize}
	\item Fehlersymptom: Daten zur Überprüfung konnten im Test nicht überprüft werden oder wurden falsch ausgewertet.
	\item Fehlerursache: Die Überprüfung verwendete einen Regex-Ausdruck, der Zeichen enthielt, die falsch interpretiert wurden.
	\item Fehlerbehebung: Zeichen wie \{ und [ wurden mit einem doppelten $\backslash$ versehen um sie erkenntlich zu machen.
\end{itemize}

\paragraph{Java Version}
\begin{itemize}
	\item Fehlersymptom: Neuere Java Funktionalitäten können nicht verwendet werden
	\item Fehlerursache: Es wurde die falsche Java-Version im Classpath verwendet
	\item Fehlerbehebung: Update auf JavaSE-1.8 mit JDK
\end{itemize}

\paragraph{Sich wiederholender Verbindungsaufbau}
\begin{itemize}
	\item Fehlersymptom: Verbindung des Graphite-Sender wird zu häufig geschlossen
	\item Fehlerursache: Bei jeder Nachricht wird das Socket überschrieben
	\item Fehlerbehebung: SocketManager Klasse wurde erstellt
\end{itemize}

\paragraph{Mehrfache nicht synchronisierte Instanzen}
\begin{itemize}
	\item Fehlersymptom: Mehrere Einheiten einer Klasse werden erstellt obwohl sie die gleichen Daten beinhalten müssen (multithreading)
	\item Fehlerursache: Offener Konstruktor und eigene Daten pro Instanz
	\item Fehlerbehebung: synchronisiertes Singleton
\end{itemize}

\paragraph{Grid Update-Prozess}
\begin{itemize}
	\item Fehlersymptom: Schlecht möglich die Reihenfolge von Operationen eines Grids beim Updaten einheitlich zu ändern
	\item Fehlerursache: update() wurde intern im Grid ausgeführt
	\item Fehlerbehebung: GeoGridManager erstellt, der Update-Prozess zentralisiert
\end{itemize}
	\include{./testTabellen/testTabellen}

\end{document}